
\section{劈尖薄膜的干涉}

\begin{conclusion}[劈尖薄膜的干涉]
两块平板夹成很小夹角时，中间的空气层或透明薄层厚度沿某一方向近似线性变化，这种薄膜称为\textbf{劈尖}。设劈尖角为 \(\alpha\)，薄膜折射率为 \(n\)，入射单色光波长为 \(\lambda\)。从棱边起沿垂直棱边方向量距离 \(x\)，则
\[
d=x\sin\alpha\simeq x\alpha\qquad(\alpha\ll 1).
\]
厚度相同的点光程差相同，所以劈尖的等厚条纹是一组与棱边平行的明暗相间直条纹。

反射光中若两束相干光之间只有一次半波损失，则
\[
\delta=2nd+\frac{\lambda}{2}.
\]
明暗条件为
\[
\begin{cases}
2nd+\dfrac{\lambda}{2}=k\lambda, & \text{明纹},\quad k=1,2,3,\dots,\\[6pt]
2nd+\dfrac{\lambda}{2}=(2k+1)\dfrac{\lambda}{2}, & \text{暗纹},\quad k=0,1,2,\dots.
\end{cases}
\]
因此棱边处 \(d=0\) 为暗纹；第 \(k\) 级明纹和第 \(k\) 级暗纹对应的厚度分别为
\[
d_k^{(\text{明})}=\frac{2k-1}{4n}\lambda\quad(k=1,2,3,\dots),
\qquad
d_k^{(\text{暗})}=\frac{k}{2n}\lambda\quad(k=0,1,2,\dots).
\]
相邻两条明纹或相邻两条暗纹的厚度差相同：
\[
\Delta d=d_{k+1}-d_k=\frac{\lambda}{2n}.
\]
若相邻同类条纹在表面上的距离为 \(l\)，则
\[
l\sin\alpha=\Delta d=\frac{\lambda}{2n},
\qquad
l=\frac{\lambda}{2n\sin\alpha}\simeq \frac{\lambda}{2n\alpha}.
\]
波长一定时，劈尖角越大，条纹间距越小，条纹越密；劈尖角过大时，条纹会难以分辨。

若两次反射的半波损失数相同，则附加的半波损失相互抵消，光程差改为
\[
\delta=2nd.
\]
这时棱边 \(d=0\) 为明纹，暗纹条件为
\[
2nd=(2k+1)\frac{\lambda}{2}\qquad(k=0,1,2,\dots).
\]
例如薄膜折射率为 \(n_2\)，且 \(n_1<n_2<n_3\) 时，上、下表面的反射都发生半波损失。若从棱边 \(O\) 起向右数，第五条暗纹落在薄膜表面某处，则 \(k=4\)，于是
\[
2n_2d=\frac{9}{2}\lambda,
\qquad
d=\frac{9\lambda}{4n_2}
=4.5\,\frac{\lambda}{2n_2}.
\]

劈尖干涉常用于测量小厚度和小形变：
\begin{enumerate}
  \item \textbf{测镀膜厚度。} 将基板上的镀膜腐蚀成劈尖形，用单色光垂直照射。由条纹级次或相邻条纹的厚度差 \(\lambda/(2n)\)，即可反推出薄膜厚度。
  \item \textbf{测细丝直径。} 把金属丝夹在两块平玻璃之间形成空气劈尖。若棱边到金属丝的距离为 \(L\)，相邻同类条纹间距为 \(l\)，金属丝直径为 \(D\)，则
  \[
  \alpha\simeq \frac{D}{L},
  \qquad
  l=\frac{\lambda}{2n\alpha},
  \qquad
  D=L\alpha=\frac{L\lambda}{2nl}.
  \]
  对空气劈尖取 \(n\simeq1\)。若 \(\lambda=\SI{589.3}{nm}\)，\(L=\SI{28.880}{mm}\)，30 条明纹间距总长为 \(\SI{4.295}{mm}\)，则
  \[
  l=\frac{4.295}{29}\,\si{mm},
  \qquad
  D=\frac{28.880\times 589.3\times 10^{-6}}{2(4.295/29)}\,\si{mm}
  \approx \SI{5.74e-2}{mm}.
  \]
  \item \textbf{测微小厚度改变量。} 当 \(\alpha\) 不变时，整体增加薄膜厚度不会改变条纹间距，只会使条纹整体移动。对一次半波损失的反射劈尖，
  \[
  2nd_k+\frac{\lambda}{2}=k\lambda
  \quad\Longrightarrow\quad
  k=\frac{2n}{\lambda}d_k+\frac{1}{2}.
  \]
  若厚度增加 \(\lambda/(2n)\)，则 \(k\to k+1\)，即高一级条纹移到原来低一级条纹的位置，整个干涉图样向干涉级低的方向移动，也就是向棱边移动。若条纹平移距离为 \(A\)，则
  \[
  \frac{A}{l}=\frac{H}{\lambda/(2n)}
  \qquad\Longrightarrow\qquad
  H=\frac{\lambda A}{2nl}.
  \]
  若用条纹移动测热膨胀，长度由 \(l_0\) 变到 \(l\)，温度由 \(t_0\) 变到 \(t\)，且条纹移动 \(N\) 条，则
  \[
  \Delta l=l-l_0=N\frac{\lambda}{2},
  \qquad
  \beta=\frac{\Delta l}{l_0(t-t_0)}
  =\frac{N\lambda}{2l_0(t-t_0)}.
  \]
  \item \textbf{检查光学表面质量。} 将待检表面 \(B\) 与参考平面组成空气劈尖。若 \(B\) 平整，条纹应为平直等距条纹；若 \(B\) 有凹坑或凸起，局部厚度改变会使条纹弯曲。设相邻条纹间距为 \(b\)，局部条纹偏移为 \(a\)，缺陷高度为 \(H\)，则
  \[
  \frac{a}{b}=\frac{H}{\lambda/2}
  \qquad\Longrightarrow\qquad
  H=\frac{\lambda}{2}\frac{a}{b}
  \]
  对空气劈尖成立；一般介质中把 \(\lambda/2\) 改为 \(\lambda/(2n)\)。
\end{enumerate}
\end{conclusion}

\begin{exercise}[2014级期末：滚柱间距变小时劈尖条纹数和间距怎样变]
两根直径有微小差别、彼此平行的滚柱之间距离为 \(L\)，其上放两块平板玻璃，中间形成空气劈尖。单色光垂直入射时产生等厚干涉条纹。若滚柱之间的距离 \(L\) 变小，判断在两滚柱之间观察到的条纹数和条纹间距怎样变化。
\end{exercise}

\begin{solution}
两滚柱直径差给定，所以两端的空气膜厚度差 \(\Delta d\) 不变；改变 \(L\) 只是改变劈尖角。

劈尖角小量近似为 \(\alpha\approx\Delta d/L\)。相邻同类条纹间距和两滚柱之间的总条纹数分别为
 \(\Delta x=\frac{\lambda}{2\alpha} =\frac{\lambda L}{2\Delta d},\qquad N\approx\frac{\Delta d}{\lambda/2} =\frac{2\Delta d}{\lambda}.\) 
当 \(L\) 变小时，\(\Delta d\) 不变，所以条纹数 \(N\) 不变，而条纹间距 \(\Delta x\) 变小。条纹数看总厚度差，条纹间距看厚度随位置变化的斜率；把两者混在一起，就会误以为间距变小必然让条纹数增加。
\end{solution}

\begin{exercise}[2006级期末：空气劈形膜上玻璃平移时条纹怎样动]
两块平玻璃构成空气劈形膜，左边为棱边，用单色平行光垂直入射。若上面的平玻璃慢慢地向上平移，判断干涉条纹的移动方向和条纹间隔变化。
\end{exercise}

\begin{solution}
劈尖条纹是等厚线。上玻璃若整体向上平移，劈尖角不变，因此厚度随位置的斜率不变；但每个位置的空气膜厚度都会增加同一个常量。

设平移前空气膜厚度近似为 \(d(x)=\alpha x\)，上板整体上移 \(h\) 后变为 \(d'(x)=h+\alpha x\)。某一条同级条纹要求膜厚等于固定值 \(d_k\)，因此新旧位置满足
 \(h+\alpha x'=d_k,\qquad \alpha x=d_k \implies x'=x-\frac{h}{\alpha}.\) 
所以条纹向棱边方向平移，条纹间隔不变。条纹间隔由斜率 \(\alpha\) 决定；平移只给厚度加常量，不改斜率，不能把“厚度增加”误判成“劈尖角变大”。
\end{solution}


\begin{exercise}[2008级期末：空气劈尖第四条明纹随劈尖角改变的位移]
用波长为 \(\lambda\) 的单色光垂直照射由两块平玻璃板构成的空气劈形膜，已知劈尖角为 \(\theta\)。若劈尖角变为 \(\theta'\)，求从劈棱数起的第四条明条纹位移值 \(\Delta x\)。
\end{exercise}

\begin{solution}
空气劈尖在反射光中有一次半波损失。距离劈棱 \(x\) 处膜厚近似为 \(d=x\theta\)，反射明纹满足总光程差为整数倍波长。

从劈棱起数，棱边处为暗纹，第四条明纹对应
 \(2x_4\theta+\frac{\lambda}{2}=4\lambda.\) 
于是
 \(x_4=\frac{7\lambda}{4\theta},\qquad x_4'=\frac{7\lambda}{4\theta'}.\) 
位移按新位置减旧位置记：
 \(\Delta x=x_4'-x_4 =\frac{7\lambda}{4}\left(\frac{1}{\theta'}-\frac{1}{\theta}\right) =\frac{7\lambda(\theta-\theta')}{4\theta\theta'}.\) 
第四条明纹离棱边不是 \(4\) 个条纹间距，而是 \(3.5\) 个条纹间距，因为反射光在棱边处为暗纹。若 \(\theta'\) 变小，条纹间距变大，第四明纹应向远离棱边方向移动，公式也给出正位移。
\end{solution}


\begin{exercise}[用劈尖条纹间距测金属细丝直径时为什么先写几何斜率]
把金属细丝夹在两块平玻璃之间形成空气劈尖。用波长 \(\lambda=\SI{589.3}{nm}\) 的单色光垂直照射，已知金属丝到劈尖顶点的距离 \(L=\SI{28.880}{mm}\)，测得 \(30\) 条明纹跨越的总长度为 \(\SI{4.295}{mm}\)。求金属细丝直径 \(D\)。
\end{exercise}

\begin{solution}
这类题看上去在测长度，其实真正要写的是“厚度沿长度方向怎样线性增长”。也就是说，先把劈尖的几何斜率写成 \(\alpha\approx D/L\)，再用条纹间距把 \(\alpha\) 反推出去。

这里空气劈尖取 \(n=1\)，且 \(30\) 条明纹跨越的是 \(29\) 个条纹间距，所以
\[
l=\frac{4.295}{29}\,\si{mm}\approx \SI{0.1481}{mm},\qquad
D=L\alpha=L\frac{\lambda}{2l}
=\frac{28.880\times 589.3\times10^{-6}}{2\times0.1481}\,\si{mm}
\approx \SI{0.0574}{mm}.
\]
细丝直径约为 \(\SI{5.74e-2}{mm}=\SI{57.4}{\micro m}\)，落在几十微米量级是合理的。这里不能把 \(30\) 条条纹误当成 \(30\) 个间距，也不能漏掉空气层对应 \(n=1\)。
\end{solution}


\begin{exercise}[空气劈尖条纹弯向厚端时怎样判断凸凹缺陷]
光学平板玻璃 \(A\) 与待测工件 \(B\) 间形成空气劈尖，左侧为空气膜薄端、右侧为空气膜厚端。用波长 \(\lambda=\SI{500}{nm}\) 的单色光垂直照射，反射光干涉条纹中弯曲部分向厚端偏移，且顶点恰好与其右边相邻直条纹的连线相切。判断工件上表面缺陷是凸起还是凹槽，并求最大高度或深度。
\end{exercise}

\begin{solution}
题目说弯曲顶点恰好与右边相邻直条纹相切，表示缺陷导致的最大膜厚变化正好相当于相邻一级条纹，即
 \(h_{\max}=\frac{\lambda}{2} =\frac{\SI{500}{nm}}{2} =\SI{250}{nm}.\) 
所以缺陷为凸起，最大高度为 \(\SI{250}{nm}\)。
\end{solution}
